\chapter{Programs — Simulation}
%%% generated by the blueprint pipeline from TCSlib.Programs.Light.Compiler.Simulation
\section{Overview}

Declarations of the Lean module \texttt{TCSlib.Programs.Light.Compiler.Simulation}.

\section{Declarations}

\begin{theorem}[Instruction locations in an assignment]
\lean{Light.Compiler.codeAt_set}
\label{Light.Compiler.codeAt_set}
\leanok
\uses{Light.Compiler.cADDR, Light.Compiler.cFP, Light.Compiler.cPool, Light.Compiler.cT, Light.Compiler.compileExpr, Light.Compiler.compileStmt, Light.Compiler.length_compileExpr, Light.Expr.size, WordRAM.CodeAt, WordRAM.CodeAt.cast_pos, WordRAM.CodeAt.left, WordRAM.CodeAt.right}
\difficulty{1}
Let `code` be a WordRAM instruction list, $G$ a code layout, $pos$ a natural number, $x$
a variable index, and $e$ an expression. If the compiled assignment of $e$ to variable
$x$ at $pos$ is a prefix of the code beginning at $pos$, then the compiled expression is
a prefix beginning at $pos$, and the two-instruction sequence that computes the
variable's address and stores the expression value is a prefix beginning at $pos+|e|$.
\end{theorem}

\begin{theorem}[Instruction locations in a sequence]
\lean{Light.Compiler.codeAt_seq}
\label{Light.Compiler.codeAt_seq}
\leanok
\uses{Light.Compiler.compileStmt, Light.Compiler.length_compileStmt, Light.Stmt.size, WordRAM.CodeAt, WordRAM.CodeAt.cast_pos, WordRAM.CodeAt.left, WordRAM.CodeAt.right}
\difficulty{1}
Let `code` be a WordRAM instruction list, $G$ a code layout, $pos$ a natural number, and
$s,t$ statements. If the compilation of the sequential composition of $s$ followed by
$t$ at $pos$ is a prefix of the code beginning at $pos$, then the compiled code for $s$
is a prefix beginning at $pos$, and the compiled code for $t$ is a prefix beginning at
$pos+|s|_G$, where $|s|_G$ is the size of $s$ for frame size $G.F$.
\end{theorem}

\begin{theorem}[Instruction locations in a conditional]
\lean{Light.Compiler.codeAt_ite}
\label{Light.Compiler.codeAt_ite}
\leanok
\uses{Light.Compiler.cNEG, Light.Compiler.compileCond, Light.Compiler.compileStmt, Light.Compiler.length_compileCond, Light.Compiler.length_compileStmt, Light.Cond.size, Light.Stmt.size, WordRAM.CodeAt, WordRAM.CodeAt.cast_pos, WordRAM.CodeAt.left, WordRAM.CodeAt.right}
\difficulty{3}
Let $P$ be a list of WordRAM instructions, let $G$ be a code layout, let $q$ be a
natural-number position, let $c$ be a condition, and let $s$ and $t$ be statements. If
the code obtained by compiling “if $c$ then $s$ else $t$” at position $q$ is a prefix of
the suffix of $P$ beginning at $q$, then the code for $c$ is a prefix of the suffix
beginning at $q$; the code for $s$ is a prefix of the suffix beginning at $q+|c|$; the
instruction that jumps to $q+|c|+|s|_G+1+|t|_G$ is at position $q+|c|+|s|_G$; and the
code for $t$ is a prefix of the suffix beginning at $q+|c|+|s|_G+1$. Here $|c|$ is the
condition size, and $|s|_G,|t|_G$ are the statement sizes for frame size $G.F$.
\end{theorem}

\begin{theorem}[Instruction locations in a while loop]
\lean{Light.Compiler.codeAt_while}
\label{Light.Compiler.codeAt_while}
\leanok
\uses{Light.Compiler.cNEG, Light.Compiler.compileCond, Light.Compiler.compileStmt, Light.Compiler.length_compileCond, Light.Compiler.length_compileStmt, Light.Cond.size, Light.Stmt.size, WordRAM.CodeAt, WordRAM.CodeAt.cast_pos, WordRAM.CodeAt.left, WordRAM.CodeAt.right}
\difficulty{3}
Let $P$ be a finite list of WordRAM instructions, let $G$ be a code layout, let $pos$ be
a natural number, let $c$ be a condition, and let $s$ be a statement. If the instruction
list produced by compiling the while loop with condition $c$ and body $s$ at position
$pos$ is a prefix of the suffix of $P$ beginning at $pos$, then the compiled condition
is a prefix of $P$ beginning at $pos$, the compiled body is a prefix beginning at $pos +
\mathrm{size}(c)$, and the single instruction that conditionally jumps back to $pos$ is
a prefix beginning at $pos + \mathrm{size}(c) + \mathrm{size}_{G.F}(s)$.
\end{theorem}

\begin{theorem}[Instruction segments for a compiled call]
\lean{Light.Compiler.codeAt_call}
\label{Light.Compiler.codeAt_call}
\leanok
\uses{Light.Compiler.argsCode, Light.Compiler.argsSize, Light.Compiler.cADDR, Light.Compiler.cFP, Light.Compiler.cNEG, Light.Compiler.cPool, Light.Compiler.cRET, Light.Compiler.callCode, Light.Compiler.callSize, Light.Compiler.compileStmt, Light.Compiler.length_argsCode, Light.Compiler.length_newFrame, Light.Compiler.newFrame, WordRAM.CodeAt, WordRAM.CodeAt.cast_pos, WordRAM.CodeAt.left, WordRAM.CodeAt.right}
\difficulty{3}
Let $G$ be a code layout with $F$ local variables per frame and a list of procedure
entry positions, and let $code$ be a finite instruction list. Suppose the compiled call
to procedure $p$ with arguments $args$, storing its result in variable $x$ and starting
at position $pos$, is a prefix of the code beginning at $pos$. Then the
argument-evaluation instructions followed by the instructions that initialise a new
frame form a prefix beginning at $pos$; the branch to procedure $p$’s entry position (or
position $0$ if the list has no entry at index $p$) forms a prefix beginning at
$pos+\mathrm{argsSize}(args)+2F+4$; and the two instructions that store the call result
form a prefix beginning at $pos+\mathrm{callSize}(F,args)$.
\end{theorem}

\begin{definition}[Compiler setting]
\lean{Light.Compiler.Setting}
\label{Light.Compiler.Setting}
\leanok
\uses{Light.Compiler.Sizes, Light.Compiler.compileStmt, Light.Compiler.dispatcher, Light.Compiler.epilogue, Light.Compiler.epilogueSize, Light.Program, Light.Stmt, Light.Stmt.size, Light.Stmt.width, WordRAM.CodeAt, WordRAM.Instr}
A compiler setting for word size $W$ consists of size and layout data, a light program
$P$, and a machine-code list. It includes a local-variable bound $F$, procedure entry
positions, a dispatcher position, run limits, and a stack-cell count $Q$; the word size
must satisfy the stated fit inequalities for the limits, stack, and dispatcher-sized
pool, and the pool bound must satisfy $2(F+1)\leq\mathrm{disp}$. Every procedure body in
$P$ must have compiled code at its entry followed by its return sequence, with the
resulting code ending no later than the dispatcher, and every body must have width at
most $F$; the dispatcher code must occur at the dispatcher position.
\end{definition}

\begin{definition}[Machine steps per language step]
\lean{Light.Compiler.stepsPerStep}
\label{Light.Compiler.stepsPerStep}
\leanok
For each natural number $n$, the number of machine steps assigned per language step is
$2n$.
\end{definition}

\begin{theorem}[Multiplicative step bound]
\lean{Light.Compiler.le_stepsPerStep_mul}
\label{Light.Compiler.le_stepsPerStep_mul}
\leanok
\uses{Light.Compiler.stepsPerStep}
\difficulty{3}
For natural numbers $a$, $c$, and $n$, if $a\le n$ and $1\le c$, then $a\le (2n)c$.
\end{theorem}

\begin{theorem}[Addition bound for step counts]
\lean{Light.Compiler.add_le_stepsPerStep_mul}
\label{Light.Compiler.add_le_stepsPerStep_mul}
\leanok
\uses{Light.Compiler.le_stepsPerStep_mul, Light.Compiler.stepsPerStep}
\difficulty{1}
For natural numbers $a,b,c,k,n$, if $a\le n$, $1\le c$, and $b\le 2nk$, then $a+b\le
2n(c+k)$.
\end{theorem}

\begin{theorem}[Bound for Call-Step Cost]
\lean{Light.Compiler.call_le_stepsPerStep_mul}
\label{Light.Compiler.call_le_stepsPerStep_mul}
\leanok
\uses{Light.Compiler.stepsPerStep}
\difficulty{3}
For natural numbers $pos,size,b,r,a,k,n$, suppose $pos+(size+2)\le n$, $r\le n$, and
$b\le 2nk$. Then $size+b+r+2(pos+size+1)+2\le 2n(a+2+k)$.
\end{theorem}

\begin{definition}[Starting configuration for a statement]
\lean{Light.Compiler.Start}
\label{Light.Compiler.Start}
\leanok
\uses{Light.Compiler.Framed, Light.Compiler.compileStmt, Light.State, Light.State.Bounded, Light.Stmt, Light.Stmt.size, Light.Stmt.width, WordRAM.CodeAt}
Let $W$ be a word size and let $S$ be a compiler setting with frame size $F$, dispatcher
position $D$, stack capacity $Q$, limits $\text{lim}$, code, and code layout. For a
statement $s$, state $σ$, code position $p$, stack-frame index $q$, and machine memory
map $m:\mathbb{Z}\to\mathrm{BitVec}(W)$, suppose the distinguished memory locations
encode $1$, $0$, and $-1$ at addresses $-5$, $-7$, and $-9$, respectively; for every
$n\le D$, address $-(25+4n)$ encodes $n$; address $-11$ encodes the stack pointer $-2q$;
each frame-local address $-2(q+x)$ encodes $σ.loc(x)$ for $x<F$; and each address $a$
encodes $σ.mem(a)$ for $a<\text{lim.space}$, with all integers represented as $W$-bit
words. Assume also that $1\le q$, $q+F\le Q$, the compiled code for $s$ in the setting's
code layout is located at $p$ and its size fits before $D$, the width of $s$ is at most
$F$, and $σ$ is bounded by the limits. Then the stack has room for the current frame
and for $F+1$ additional cells at each of the remaining $\text{lim.depth}-d$ call
levels: $q+F+(F+1)(\text{lim.depth}-d)\le Q$.
\end{definition}

\begin{definition}[Completed execution of a statement]
\lean{Light.Compiler.Finished}
\label{Light.Compiler.Finished}
\leanok
\uses{Light.Compiler.Rel, Light.Compiler.Writable, Light.Compiler.stepsPerStep, Light.State, Light.Stmt, Light.Stmt.size, WordRAM.AgreeOutside, WordRAM.Steps}
Fix a word size $W$, compiler setting $S$, statement $s$, resulting state $\sigma'$,
natural numbers $c,\mathrm{pos},q$, and memory maps $m,m'$. A completed execution
consists of a run of the compiled code from position $\mathrm{pos}$ to the position
immediately after its code, taking $n$ machine steps with $n\leq 2\,\mathrm{disp}\,c$;
the final memory represents $\sigma'$ in the same frame, and cells outside those
writable by a statement in frame $q$ are unchanged.
\end{definition}

\begin{definition}[Statement simulation]
\lean{Light.Compiler.Simulates}
\label{Light.Compiler.Simulates}
\leanok
\uses{Light.Compiler.Finished, Light.Compiler.Start, Light.State, Light.Stmt}
Fix a word size $W$, a compiler setting $S$, a natural number $d$, a statement $s$,
states $\sigma$ and $\sigma'$, and a natural number $c$. The statement $s$ is simulated
from $\sigma$ to $\sigma'$ with parameters $d$ and $c$ if, for every code position
$pos$, frame index $q$, and memory map $m : \mathbb{Z} \to \mathrm{BitVec}(W)$ for which
$\mathrm{Start}(S,d,s,\sigma,pos,q,m)$ holds, there are a memory map $m' : \mathbb{Z}
\to \mathrm{BitVec}(W)$ and a natural number $n$ such that
$\mathrm{Finished}(S,s,\sigma',c,pos,q,m,m',n)$ holds.
\end{definition}

\begin{theorem}[Reestablishing a start configuration]
\lean{Light.Compiler.Start.sub}
\label{Light.Compiler.Start.sub}
\leanok
\uses{Light.Compiler.Rel, Light.Compiler.Start, Light.Compiler.compileStmt, Light.State.Bounded, Light.Stmt, Light.Stmt.size, Light.Stmt.width, WordRAM.CodeAt}
\difficulty{1}
Let $S$ be a setting, and suppose a start configuration for a statement $s$ at depth
$d$, code position $\mathrm{pos}$, and stack position $q$ satisfies the start
conditions. Let $t$ be a statement with compiled code at position $\mathrm{pos}'$, whose
code length ends at or before the dispatcher, and whose width is at most $S.F$. Suppose
$\sigma'$ is a state bounded by the limits of $S$, and a word-RAM memory $m'$ is related
to $\sigma'$ at stack position $q$. Then $t$, $\sigma'$, $\mathrm{pos}'$, and $m'$
satisfy the start conditions for $S$ at depth $d$ and stack position $q$.
\end{theorem}

\begin{theorem}[Stack room at the next depth]
\lean{Light.Compiler.Start.room_succ}
\label{Light.Compiler.Start.room_succ}
\leanok
\uses{Light.Compiler.Start}
\difficulty{3}
Let a compiler setting have frame size $F$, stack capacity $Q$, and depth limit $D$. If
a statement starts at depth $d$ and stack position $q$ in a configuration satisfying the
start conditions, and $d<D$, then $q+2F+1+(F+1)(D-(d+1))\le Q$.
\end{theorem}

\begin{theorem}[Compiled conditions decide their truth]
\lean{Light.Compiler.Start.test}
\label{Light.Compiler.Start.test}
\leanok
\uses{Light.Compiler.Decides, Light.Compiler.Start, Light.Compiler.compileCond, Light.Compiler.decides_compileCond, Light.Cond, Light.Cond.Holds, Light.Cond.Safe, Light.Cond.size, Light.Cond.width, WordRAM.CodeAt}
\difficulty{1}
Let $S$ be a setting, and suppose a start configuration for a statement $s$ at depth
$d$, code position $\mathrm{pos}$, and stack position $q$ satisfies the start
conditions, with word-RAM memory $m$. Let $c$ be a condition whose compiled code occurs
at $\mathrm{pos}$, whose width is at most $S.F$, and which is safe for a light state
$\sigma$ under the limits of $S$. Then there exist a word-RAM memory $m'$ and a natural
number $k\le\mathrm{size}(c)$ such that $m'$ agrees with $m$ outside the scratch-address
set for frame size $S.F$; if $c$ holds in $\sigma$, the compiled code takes $k$ steps
from position $\mathrm{pos}$ to $\mathrm{pos}+\mathrm{size}(c)$, and if $c$ does not
hold, it takes $k$ steps from $\mathrm{pos}$ to branch position $l$.
\end{theorem}

\begin{theorem}[Simulation of the Skip Statement]
\lean{Light.Compiler.sim_skip}
\label{Light.Compiler.sim_skip}
\leanok
\uses{Light.Compiler.Simulates, WordRAM.AgreeOutside.refl, WordRAM.Steps.refl}
\difficulty{1}
For every compiler setting $S$, natural number $d$, and state $\sigma$, the skip
statement simulates execution from $\sigma$ to itself with cost zero. For every
admissible initial machine placement and memory encoding, the compiled code reaches its
endpoint within the corresponding machine-step bound, with the memory still encoding
$\sigma$ and unchanged outside writable addresses.
\end{theorem}

\begin{theorem}[Simulation of Variable Assignment]
\lean{Light.Compiler.sim_set}
\label{Light.Compiler.sim_set}
\leanok
\uses{Light.Compiler.Framed.same, Light.Compiler.Simulates, Light.Compiler.codeAt_set, Light.Compiler.effects_compileExpr_of_width, Light.Compiler.le_stepsPerStep_mul, Light.Compiler.length_compileExpr, Light.Compiler.scratch_subset_writable, Light.Compiler.steps_setVar, Light.Compiler.straight_compileExpr, Light.Compiler.temps_subset_scratch, Light.Expr, Light.Expr.Safe, Light.Expr.cost, Light.Expr.size, Light.Expr.val, Light.Stmt.size, Light.Stmt.width_set_le_iff, WordRAM.AgreeOutside.mono, WordRAM.AgreeOutside.trans, WordRAM.Steps.cast_pos, WordRAM.Steps.trans, WordRAM.steps_straight}
\difficulty{4}
Let $S$ be a compiler setting, let $d,x$ be natural numbers, let $e$ be an expression
safe in state $\sigma$, and let $\sigma$ be a state. Assigning the value of $e$ in
$\sigma$ to variable $x$ simulates execution from $\sigma$ to the state with its
variable function updated at $x$ to $e(\sigma)$ and with memory unchanged, in
$\mathrm{cost}(e)+1$ steps. This holds for every admissible initial machine placement
and memory encoding; the compiled code reaches its endpoint within the corresponding
machine-step bound, with memory encoding the updated state and agreement outside
writable addresses.
\end{theorem}

\begin{theorem}[Simulation of a Memory Store]
\lean{Light.Compiler.sim_store}
\label{Light.Compiler.sim_store}
\leanok
\uses{Light.Compiler.Rel.same, Light.Compiler.Rel.setMem, Light.Compiler.Simulates, Light.Compiler.Sizes.inRange_addr, Light.Compiler.codeAt_after_pair, Light.Compiler.le_stepsPerStep_mul, Light.Compiler.mem_writable_of_addr, Light.Compiler.scratch_subset_writable, Light.Compiler.steps_pair, Light.Expr, Light.Expr.Safe, Light.Expr.cost, Light.Expr.size, Light.Expr.val, Light.Limits.Addr, Light.Stmt.size, Light.Stmt.width_store_le_iff, WordRAM.AgreeOutside.mono, WordRAM.AgreeOutside.update, WordRAM.Steps.cast_pos, WordRAM.Steps.trans, WordRAM.steps_store}
\difficulty{4}
Let $S$ be a compiler setting, let $d$ be a natural number, and let $a,e$ be expressions
safe in state $\sigma$. Suppose the value of $a$ in $\sigma$ is an address within the
space limit of $S$. Then storing the value of $e$ at that address simulates execution
from $\sigma$ to the state whose memory is updated at the natural-number address
represented by $a(\sigma)$ to the value $e(\sigma)$, with variable values unchanged, in
$\mathrm{cost}(a)+\mathrm{cost}(e)+1$ steps. This holds for every admissible initial
machine placement and memory encoding; the compiled code reaches its endpoint within the
corresponding machine-step bound, with memory encoding the updated state and agreement
outside writable addresses.
\end{theorem}

\begin{theorem}[Simulation of Sequential Statements]
\lean{Light.Compiler.sim_seq}
\label{Light.Compiler.sim_seq}
\leanok
\uses{Light.Compiler.Simulates, Light.Compiler.Start.sub, Light.Compiler.codeAt_seq, Light.Exec, Light.Exec.bounded, Light.Stmt.size, Light.Stmt.width_seq_le_iff, WordRAM.AgreeOutside.trans, WordRAM.Steps.cast_pos, WordRAM.Steps.trans}
\difficulty{4}
Let $S$ be a program setting. Suppose the first statement $s_1$ executes from $\sigma$
to $\sigma_1$ in $c_1$ steps, and both $s_1$ from $\sigma$ to $\sigma_1$ in $c_1$ steps
and $s_2$ from $\sigma_1$ to $\sigma_2$ in $c_2$ steps simulate their executions. Then
the sequence of $s_1$ followed by $s_2$ simulates execution from $\sigma$ to $\sigma_2$
in $c_1+c_2$ steps.
\end{theorem}

\begin{theorem}[Simulation of the True Conditional Branch]
\lean{Light.Compiler.sim_iteTrue}
\label{Light.Compiler.sim_iteTrue}
\leanok
\uses{Light.Compiler.Decides.of_holds, Light.Compiler.Rel.same, Light.Compiler.Simulates, Light.Compiler.Sizes.inRange_neg_one, Light.Compiler.Start.sub, Light.Compiler.Start.test, Light.Compiler.add_le_stepsPerStep_mul, Light.Compiler.codeAt_ite, Light.Compiler.scratch_subset_writable, Light.Compiler.steps_jump, Light.Cond, Light.Cond.Holds, Light.Cond.Safe, Light.Cond.cost, Light.Cond.size, Light.Stmt.size, Light.Stmt.width_ite_le_iff, WordRAM.AgreeOutside.mono, WordRAM.AgreeOutside.trans, WordRAM.Steps.cast_count, WordRAM.Steps.cast_pos, WordRAM.Steps.trans}
\difficulty{5}
Let $S$ be a compiler setting, let $d$ be a natural number, and let $c$ be a condition
that is safe and true in state $\sigma$. If the true branch $s_1$ simulates execution
from $\sigma$ to $\sigma'$ in $k$ steps, then the conditional with branches $s_1$ and
$s_2$ simulates execution from $\sigma$ to $\sigma'$ in $\mathrm{cost}(c)+1+k$ steps.
Simulation here quantifies over every admissible initial machine placement and memory
encoding and guarantees a run to the code endpoint within the corresponding machine-step
bound, with memory encoding $\sigma'$ and agreement outside the writable addresses.
\end{theorem}

\begin{theorem}[Simulation of the False Conditional Branch]
\lean{Light.Compiler.sim_iteFalse}
\label{Light.Compiler.sim_iteFalse}
\leanok
\uses{Light.Compiler.Decides.of_fails, Light.Compiler.Rel.same, Light.Compiler.Simulates, Light.Compiler.Start.sub, Light.Compiler.Start.test, Light.Compiler.add_le_stepsPerStep_mul, Light.Compiler.codeAt_ite, Light.Compiler.scratch_subset_writable, Light.Cond, Light.Cond.Holds, Light.Cond.Safe, Light.Cond.cost, Light.Cond.size, Light.Stmt.size, Light.Stmt.width_ite_le_iff, WordRAM.AgreeOutside.mono, WordRAM.AgreeOutside.trans, WordRAM.Steps.cast_pos, WordRAM.Steps.trans}
\difficulty{4}
Let $S$ be a compiler setting, let $d$ be a natural number, and let $c$ be a condition
that is safe and false in state $\sigma$. If the false branch $s_2$ simulates execution
from $\sigma$ to $\sigma'$ in $k$ steps, then the conditional with branches $s_1$ and
$s_2$ simulates execution from $\sigma$ to $\sigma'$ in $\mathrm{cost}(c)+1+k$ steps.
Simulation here quantifies over every admissible initial machine placement and memory
encoding and guarantees a run to the code endpoint within the corresponding machine-step
bound, with memory encoding $\sigma'$ and agreement outside the writable addresses.
\end{theorem}

\begin{theorem}[Simulation of a False While Loop]
\lean{Light.Compiler.sim_whileFalse}
\label{Light.Compiler.sim_whileFalse}
\leanok
\uses{Light.Compiler.Decides.of_fails, Light.Compiler.Rel.same, Light.Compiler.Simulates, Light.Compiler.Start.test, Light.Compiler.codeAt_while, Light.Compiler.le_stepsPerStep_mul, Light.Compiler.scratch_subset_writable, Light.Cond, Light.Cond.Holds, Light.Cond.Safe, Light.Cond.cost, Light.Cond.size, Light.Stmt.size, Light.Stmt.width_while_le_iff, WordRAM.AgreeOutside.mono, WordRAM.Steps.cast_pos}
\difficulty{3}
Let $S$ be a compiler setting, let $d$ be a natural number, and let $c$ be a condition
safe in state $\sigma$ that does not hold there. Then the loop with condition $c$ and
body $s$ simulates execution from $\sigma$ to itself with cost $\mathrm{cost}(c)+1$. For
every admissible initial machine placement and memory encoding, the compiled code
reaches its endpoint within the corresponding machine-step bound, with the memory still
encoding $\sigma$ and unchanged outside writable addresses.
\end{theorem}

\begin{theorem}[Simulation of a true while loop]
\lean{Light.Compiler.sim_whileTrue}
\label{Light.Compiler.sim_whileTrue}
\leanok
\uses{Light.Compiler.Decides.of_holds, Light.Compiler.Rel.same, Light.Compiler.Simulates, Light.Compiler.Sizes.inRange_neg_one, Light.Compiler.Start.sub, Light.Compiler.Start.test, Light.Compiler.add_le_stepsPerStep_mul, Light.Compiler.codeAt_while, Light.Compiler.scratch_subset_writable, Light.Compiler.steps_jump, Light.Cond, Light.Cond.Holds, Light.Cond.Safe, Light.Cond.cost, Light.Cond.size, Light.Exec, Light.Exec.bounded, Light.Stmt.size, Light.Stmt.width_while_le_iff, WordRAM.AgreeOutside.mono, WordRAM.AgreeOutside.trans, WordRAM.Steps.cast_count, WordRAM.Steps.trans}
\difficulty{5}
Let $S$ be a compiler setting, let $d,k_1,k_2$ be nonnegative integers, and let $c$ be a
condition, $s$ a statement, and $\sigma,\sigma_1,\sigma_2$ states. If $c$ is safe in
$\sigma$ under the limits of $S$ and holds in $\sigma$, if executing $s$ from $\sigma$
to $\sigma_1$ at depth $d$ has cost $k_1$ and is simulated by $S$ with that cost, and if
the loop that tests $c$ and executes $s$ is simulated by $S$ from $\sigma_1$ to
$\sigma_2$ with cost $k_2$, then that loop is simulated by $S$ from $\sigma$ to
$\sigma_2$ with cost $\operatorname{cost}(c)+1+k_1+k_2$.
\end{theorem}

\begin{theorem}[Procedure call entry]
\lean{Light.Compiler.steps_callEnter}
\label{Light.Compiler.steps_callEnter}
\leanok
\uses{Light.Compiler.Framed.same, Light.Compiler.Rel, Light.Compiler.Sizes.inRange_neg_one, Light.Compiler.Start, Light.Compiler.Start.room_succ, Light.Compiler.Writable, Light.Compiler.cStack, Light.Compiler.callSize, Light.Compiler.codeAt_call, Light.Compiler.effects_argsCode, Light.Compiler.effects_newFrame, Light.Compiler.length_argsCode, Light.Compiler.length_newFrame, Light.Compiler.scratch_subset_writable, Light.Compiler.steps_jump, Light.Compiler.straight_argsCode, Light.Compiler.straight_newFrame, Light.Compiler.temps_subset_scratch, Light.Expr.Safe, Light.Expr.val, Light.Stmt.width_call_le_iff, Light.frame, WordRAM.AgreeOutside, WordRAM.AgreeOutside.mono, WordRAM.AgreeOutside.trans, WordRAM.Steps, WordRAM.Steps.cast_count, WordRAM.Steps.trans, WordRAM.effects_append, WordRAM.steps_straight, WordRAM.wd}
\difficulty{4}
Let $S$ be a compiler setting, and suppose a call statement with procedure index $p$,
argument expressions $args$, and return-variable index $x$ starts at depth $d$, code
position $pos$, stack base $q$, state $\sigma$, and word memory $m$. Assume the starting
conditions hold, every argument expression is safe in $\sigma$, $d$ is less than the depth
limit, and procedure $p$ has entry position $e$. Then there is a word memory $m_2$ such
that the compiled machine takes exactly the call-size number of steps from $(pos,m)$ to
position $e$ with memory $m_2$. This memory represents the callee state whose local
frame consists of the argument values and whose ordinary memory is inherited from $\sigma$;
its stack cell at index $q+F$ contains the return position $pos$ plus the call size; and
it agrees with $m$ outside the cells writable from stack base $q+F$.
\end{theorem}

\begin{theorem}[Return from a procedure call]
\lean{Light.Compiler.steps_callLeave}
\label{Light.Compiler.steps_callLeave}
\leanok
\uses{Light.Compiler.Rel, Light.Compiler.Rel.loc_zero, Light.Compiler.Rel.same, Light.Compiler.Scratch, Light.Compiler.Setting, Light.Compiler.Sizes.inRange_neg_one, Light.Compiler.Sizes.inRange_stack, Light.Compiler.cADDR, Light.Compiler.cFP, Light.Compiler.cRET, Light.Compiler.cRR, Light.Compiler.cStack, Light.Compiler.cStack_add, Light.Compiler.cStack_succ_add, Light.Compiler.epilogue, Light.Compiler.epilogueSize, Light.Compiler.steps_dispatcher, Light.Compiler.steps_jump, WordRAM.AgreeOutside, WordRAM.AgreeOutside.mono, WordRAM.AgreeOutside.read, WordRAM.AgreeOutside.refl, WordRAM.AgreeOutside.update, WordRAM.CodeAt, WordRAM.CodeAt.tail, WordRAM.Steps, WordRAM.Steps.trans, WordRAM.steps_add, WordRAM.steps_load, WordRAM.wd}
\difficulty{5}
Let a compiled program have word size $W$, and let its memory $m$ represent a
light-program state $\sigma$ at stack depth $q+F+1$. Suppose the code at position
$\mathrm{pos}$ begins with the procedure epilogue, the return-address cell at stack
position $q+F$ contains the word encoding $\mathrm{ret}$, $\mathrm{ret}<\mathrm{disp}$,
$q+F+1\le Q$, and $0<F$. Then after $5+2(\mathrm{ret}+1)$ machine steps, the program can
reach program counter $\mathrm{ret}$ with a memory $m'$ whose frame pointer contains the
word encoding the stack position $q$ and whose return register contains the word
encoding $\sigma.loc(0)$. The memories $m$ and $m'$ agree at every address outside the
union of the frame-pointer address and the scratch-address set for $F$.
\end{theorem}

\begin{theorem}[Simulation of Procedure Calls]
\lean{Light.Compiler.sim_call}
\label{Light.Compiler.sim_call}
\leanok
\uses{Light.Compiler.Framed, Light.Compiler.Rel.of_agreeOutside, Light.Compiler.Simulates, Light.Compiler.Start.room_succ, Light.Compiler.Writable, Light.Compiler.cStack, Light.Compiler.callSize, Light.Compiler.call_le_stepsPerStep_mul, Light.Compiler.codeAt_call, Light.Compiler.length_compileStmt, Light.Compiler.scratch_subset_writable, Light.Compiler.steps_callEnter, Light.Compiler.steps_callLeave, Light.Compiler.steps_setVar, Light.Compiler.writable_subset_writable, Light.Expr.Safe, Light.Expr.cost, Light.Expr.val, Light.Stmt, Light.Stmt.size, Light.Stmt.width_call_le_iff, Light.bounded_callFrame, Light.frame, WordRAM.AgreeOutside, WordRAM.AgreeOutside.mono, WordRAM.AgreeOutside.read, WordRAM.AgreeOutside.trans, WordRAM.CodeAt.cast_pos, WordRAM.CodeAt.left, WordRAM.CodeAt.right, WordRAM.Steps.cast_pos, WordRAM.Steps.trans}
\difficulty{6}
Let $S$ be a compiler setting, let $d,p,x,k$ be natural numbers, let
$\sigma,\sigma'$ be program states, and let $\mathrm{args}$ be a finite list
of expressions. Suppose every argument expression is safe in $\sigma$, procedure $p$
has body $\mathrm{body}$ in $S$, and $d$ is less than the depth limit. If
$\mathrm{body}$ simulates execution from depth $d+1$ and the state whose variable
values are the argument values (with missing arguments set to zero) and whose memory is
that of $\sigma$, to $\sigma'$ in $k$ steps, then the call to $p$ with arguments
$\mathrm{args}$, assigning its result to variable $x$, simulates execution from
$\sigma$ to the state obtained by updating variable $x$ to
$\sigma'.\mathrm{loc}(0)$ and using memory $\sigma'.\mathrm{mem}$, in
$\sum_{e\in\mathrm{args}}\mathrm{cost}(e)+2+k$ steps.
\end{theorem}

\begin{theorem}[Simulation of Executions]
\lean{Light.Compiler.sim}
\label{Light.Compiler.sim}
\leanok
\uses{Light.Compiler.Setting, Light.Compiler.Simulates, Light.Compiler.sim_call, Light.Compiler.sim_iteFalse, Light.Compiler.sim_iteTrue, Light.Compiler.sim_seq, Light.Compiler.sim_set, Light.Compiler.sim_skip, Light.Compiler.sim_store, Light.Compiler.sim_whileFalse, Light.Compiler.sim_whileTrue, Light.Exec}
\difficulty{4}
Let $S$ be a program setting, and suppose a statement $s$ executes from state $\sigma$
to state $\sigma'$ at depth $d$ in $c$ language steps. For every program position and
stack index, and every initial machine memory satisfying the starting-frame conditions
(including the code placement, bounds, state representation, and stack-space
requirements), there are a final machine memory and a number $n$ of machine steps such
that the machine reaches the position immediately after the compiled statement, $n \le
2S.\mathrm{disp}\,c$, the final memory represents $\sigma'$ in the same frame, and the
initial and final memories agree outside the writable addresses for that frame.
\end{theorem}
